\documentclass[12pt,a4paper]{article}
\usepackage[top=2.5cm,bottom=2.5cm,left=2.5cm,right=2.5cm,headheight=15pt]{geometry}
\usepackage{amsmath,amssymb}
\usepackage{tikz}
\usetikzlibrary{shapes,arrows,arrows.meta,positioning,calc,backgrounds,fit,matrix}
\usepackage{booktabs,array,longtable,tabularx,multirow}
\usepackage{xcolor}
\usepackage{enumitem}
\usepackage{fancyhdr}
\usepackage{titlesec}
\usepackage{mdframed}
\usepackage{tcolorbox}
\tcbuselibrary{skins,breakable}
\usepackage{hyperref}

%% ---- Colors ----
\definecolor{folblue}{RGB}{0,82,165}
\definecolor{lightblue}{RGB}{220,235,255}
\definecolor{darkorange}{RGB}{180,70,0}
\definecolor{lightorange}{RGB}{255,235,210}
\definecolor{darkgreen}{RGB}{20,110,20}
\definecolor{lightgreen}{RGB}{200,240,205}
\definecolor{darkred}{RGB}{150,20,20}
\definecolor{lightred}{RGB}{255,210,210}
\definecolor{lightgray}{RGB}{242,242,242}

%% ---- tcolorbox styles ----
\newtcolorbox{qbox}[2][]{enhanced,breakable,
  colback=lightorange,colframe=darkorange,arc=4pt,
  title={\textbf{#2}},fonttitle=\bfseries\color{white},
  attach boxed title to top left={yshift=-2mm,xshift=4mm},
  boxed title style={colback=darkorange,arc=3pt},#1}

\newtcolorbox{solbox}[2][]{enhanced,breakable,
  colback=lightblue,colframe=folblue,arc=4pt,
  title={\textbf{#2}},fonttitle=\bfseries\color{white},
  attach boxed title to top left={yshift=-2mm,xshift=4mm},
  boxed title style={colback=folblue,arc=3pt},#1}

\newtcolorbox{keybox}[2][]{enhanced,breakable,
  colback=lightgreen,colframe=darkgreen,arc=4pt,
  title={\textbf{#2}},fonttitle=\bfseries\color{white},
  attach boxed title to top left={yshift=-2mm,xshift=4mm},
  boxed title style={colback=darkgreen,arc=3pt},#1}

\newtcolorbox{algobox}[2][]{enhanced,breakable,
  colback=lightgray,colframe=gray!60,arc=4pt,
  title={\textbf{#2}},fonttitle=\bfseries,
  attach boxed title to top left={yshift=-2mm,xshift=4mm},
  boxed title style={colback=gray!60,arc=3pt},#1}

\newtcolorbox{warnbox}[2][]{enhanced,breakable,
  colback=lightred,colframe=darkred,arc=4pt,
  title={\textbf{#2}},fonttitle=\bfseries\color{white},
  attach boxed title to top left={yshift=-2mm,xshift=4mm},
  boxed title style={colback=darkred,arc=3pt},#1}

%% ---- Page style ----
\pagestyle{fancy}\fancyhf{}
\lhead{\color{folblue}\textbf{Game Theory --- CSE 717 InfoSec}}
\rhead{\color{folblue}Exam: 17 Jun 2026}
\cfoot{--- \thepage\ ---}
\renewcommand{\headrulewidth}{0.4pt}

%% ---- Section style ----
\titleformat{\section}{\large\bfseries\color{folblue}}{}{0em}{}[\titlerule]
\titleformat{\subsection}{\normalsize\bfseries\color{darkorange}}{}{0em}{}

\begin{document}

\begin{center}
{\LARGE\bfseries\color{folblue} CSE 717 --- Information Security}\\[4pt]
{\large\bfseries Topic 5: Game Theory \& Nash Equilibrium --- Payoff Matrix, Dominant Strategy}\\[2pt]
{\large Complete Past Paper Solutions: 2020--2024}\\[6pt]
{\normalsize Source: Lc\#3 (Game Theory in Network Security)}\\[2pt]
{\normalsize\color{darkred}\textbf{2/5 years (2023, 2024 -- both most recent), 3--6 marks. Cheap, mechanical marks once the method is drilled.}}
\end{center}

\tableofcontents

%% ==========================================
\section*{Quick Reference --- Finding Nash Equilibrium (CHEAT SHEET)}
%% ==========================================

\begin{warnbox}{2020--2022: no Game Theory question found}
This topic is NEW to the 2026 syllabus pattern --- only 2023 and 2024 (the two most recent years) ask it, both as small (3-mark) parts of a larger Section-B question. Don't over-invest, but the method below is fast and mechanical once memorized.
\end{warnbox}

\begin{algobox}{Key definitions (Lc\#3)}
\begin{itemize}[noitemsep]
  \item \textbf{Player:} an entity participating in the game (e.g., Alice, Bob, Eve, Mallory).
  \item \textbf{Action / Strategy:} a choice a player can make.
  \item \textbf{Payoff:} the gain (or loss) a player receives, depending on \emph{both} players' choices.
  \item \textbf{Payoff matrix:} a table where each cell $(x,y)$ shows (Row player's payoff, Column player's payoff) for that combination of actions.
  \item \textbf{Nash Equilibrium (NE):} a combination of strategies where \textbf{neither player can improve their own payoff by unilaterally changing their own strategy}, given the other player's choice. A game can have 0, 1, or multiple NEs.
  \item \textbf{Strictly dominant strategy:} a strategy that gives a player a \textbf{strictly higher payoff no matter what the other player does} --- i.e., it is the best choice regardless of the opponent's action.
\end{itemize}
\end{algobox}

\begin{algobox}{The 6-Step "Underline" Method to Find All NEs (Lc\#3)}
\begin{enumerate}[noitemsep]
  \item Pretend you are Player 1 (row player).
  \item Assume Player 2 (column player) picks a particular action (fix a column).
  \item Looking down that column, find Player 1's \textbf{best payoff} (highest number) and underline it.
  \item Repeat Steps 2--3 for Player 2's other action(s) (the other column(s)).
  \item Now switch roles: for each row (Player 1's action fixed), find Player 2's best payoff across that row and underline it.
  \item \textbf{Any cell where BOTH numbers are underlined is a Nash Equilibrium.}
\end{enumerate}
\textbf{Dominant strategy check:} for each player, compare their payoffs for a fixed strategy across all of the opponent's actions. If one strategy gives a strictly higher payoff in \emph{every} case, it is strictly dominant.
\end{algobox}

\begin{keybox}{Worked example from slides --- $2\times2$ payoff matrix}
\[
\begin{array}{c|cc}
 & \text{Action C} & \text{Action D} \\\hline
\text{Action A} & 10,2 & 8,3 \\
\text{Action B} & 12,4 & 10,1
\end{array}
\]
\textbf{Underline pass:} If P2 plays C, P1's best is B (12$>$10) $\to$ underline 12. If P2 plays D, P1's best is B (10$>$8) $\to$ underline 10. If P1 plays A, P2's best is D (3$>$2) $\to$ underline 3. If P1 plays B, P2's best is C (4$>$1) $\to$ underline 4.

$\Rightarrow$ Cell $(B,C)=(12,4)$ has BOTH numbers underlined $\Rightarrow$ \textbf{unique NE $=(B,C)$}. Also, P1's best response is always B (12$>$10 and 10$>$8 in both columns) $\Rightarrow$ \textbf{P1 has a strictly dominant strategy: B}.

\textbf{Second slide example} (coordination game) can have \textbf{two} NEs simultaneously, e.g. $(\text{Yes},\text{Yes})=(20,20)$ and $(\text{No},\text{No})=(10,10)$ --- both are NEs even though $(20,20)$ is more efficient; game theory alone cannot predict which occurs.
\end{keybox}

%% ==========================================
\section{2024 --- Section B Q7(a): Eve/Mallory PT vs ND payoff matrix --- find all NE, dominant strategies (3)}
%% ==========================================

\begin{qbox}{Question --- 2024 Section B Q7(a)}
Eve and Mallory are choosing which cybersecurity training to attend: Penetration Testing (PT) or Network Defense (ND). Their satisfaction (payoff) depends on both their choices; they prefer attending the same training together, but each has slightly different preferences. The payoff table shows (Eve's payoff, Mallory's payoff):
\[
\begin{array}{c|cc}
 & \text{Mallory: PT} & \text{Mallory: ND} \\\hline
\text{Eve: PT} & 4,3 & 1,2 \\
\text{Eve: ND} & 2,1 & 3,4
\end{array}
\]
\begin{enumerate}[label=(\roman*)]
  \item Identify all Nash equilibria in this scenario.
  \item Does Eve have a strictly dominant strategy? Does Mallory? \hfill [3]
\end{enumerate}
\end{qbox}

\begin{solbox}{Answer --- 2024 Section B Q7(a)}
\textbf{Step 1--2 --- if Mallory plays PT}, compare Eve's payoffs down the PT column: $4$ (Eve:PT) vs $2$ (Eve:ND). Best $=4>2 \to$ \underline{underline Eve's 4} in cell (PT,PT).

\textbf{If Mallory plays ND}, compare Eve's payoffs down the ND column: $1$ (Eve:PT) vs $3$ (Eve:ND). Best $=3>1 \to$ \underline{underline Eve's 3} in cell (ND,ND).

\textbf{If Eve plays PT}, compare Mallory's payoffs across the PT row: $3$ (Mallory:PT) vs $2$ (Mallory:ND). Best $=3>2 \to$ \underline{underline Mallory's 3} in cell (PT,PT).

\textbf{If Eve plays ND}, compare Mallory's payoffs across the ND row: $1$ (Mallory:PT) vs $4$ (Mallory:ND). Best $=4>1 \to$ \underline{underline Mallory's 4} in cell (ND,ND).

\[
\begin{array}{c|cc}
 & \text{Mallory: PT} & \text{Mallory: ND} \\\hline
\text{Eve: PT} & \underline{4},\underline{3} & 1,2 \\
\text{Eve: ND} & 2,1 & \underline{3},\underline{4}
\end{array}
\]

\textbf{(i) Nash Equilibria:} both cells $(\text{PT},\text{PT})=(4,3)$ and $(\text{ND},\text{ND})=(3,4)$ have BOTH numbers underlined.
\[
\boxed{\text{NE}_1=(\text{Eve:PT},\text{Mallory:PT})=(4,3), \qquad \text{NE}_2=(\text{Eve:ND},\text{Mallory:ND})=(3,4)}
\]
This is a \textbf{coordination game} --- both players prefer to attend the same training, but disagree on which one (Eve prefers PT, Mallory prefers ND).

\textbf{(ii) Dominant strategy check:}
\begin{itemize}[noitemsep]
  \item \textbf{Eve:} if Mallory plays PT, Eve's best is PT (4$>$2); if Mallory plays ND, Eve's best is ND (3$>$1). Eve's best choice \textbf{depends on Mallory's choice} $\Rightarrow$ \textbf{Eve has NO strictly dominant strategy.}
  \item \textbf{Mallory:} if Eve plays PT, Mallory's best is PT (3$>$2); if Eve plays ND, Mallory's best is ND (4$>$1). Mallory's best choice also \textbf{depends on Eve's choice} $\Rightarrow$ \textbf{Mallory has NO strictly dominant strategy either.}
\end{itemize}
\end{solbox}

%% ==========================================
\section{2023 --- Section B Q8(a): Alice/Bob ST\#1 vs ST\#2 payoff matrix --- NE + dominant strategy (3)}
%% ==========================================

\begin{qbox}{Question --- 2023 Section B Q8(a)}
Alice and Bob develop two new network systems, but they forgot to agree on which security strategy to follow: Strategy 1 (ST\#1) or Strategy 2 (ST\#2). Their payoffs depend on which strategy they choose and whether they both apply the same one. Alice chooses a row, Bob chooses a column, and Alice's payoff is listed first:
\[
\begin{array}{c|cc}
 & \text{Bob: ST\#1} & \text{Bob: ST\#2} \\\hline
\text{Alice: ST\#1} & 10,10 & 8,4 \\
\text{Alice: ST\#2} & 4,8 & 5,5
\end{array}
\]
\begin{enumerate}[label=(\roman*)]
  \item Find the Nash equilibrium of this scenario, if any.
  \item Does Alice have a strictly dominant strategy? Does Bob? \hfill [3]
\end{enumerate}
\end{qbox}

\begin{solbox}{Answer --- 2023 Section B Q8(a)}
\textbf{If Bob plays ST\#1}, Alice's payoffs: $10$ (ST\#1) vs $4$ (ST\#2). Best $=10>4 \to$ \underline{underline Alice's 10}.

\textbf{If Bob plays ST\#2}, Alice's payoffs: $8$ (ST\#1) vs $5$ (ST\#2). Best $=8>5 \to$ \underline{underline Alice's 8}.

\textbf{If Alice plays ST\#1}, Bob's payoffs: $10$ (ST\#1) vs $4$ (ST\#2). Best $=10>4 \to$ \underline{underline Bob's 10}.

\textbf{If Alice plays ST\#2}, Bob's payoffs: $8$ (ST\#1) vs $5$ (ST\#2). Best $=8>5 \to$ \underline{underline Bob's 8}.

\[
\begin{array}{c|cc}
 & \text{Bob: ST\#1} & \text{Bob: ST\#2} \\\hline
\text{Alice: ST\#1} & \underline{10},\underline{10} & \underline{8},4 \\
\text{Alice: ST\#2} & 4,\underline{8} & 5,5
\end{array}
\]

\textbf{(i) Nash Equilibrium:} only cell $(\text{ST\#1},\text{ST\#1})=(10,10)$ has BOTH numbers underlined.
\[
\boxed{\text{NE}=(\text{Alice:ST\#1},\text{Bob:ST\#1})=(10,10)} \quad \text{--- unique NE}
\]

\textbf{(ii) Dominant strategy check:}
\begin{itemize}[noitemsep]
  \item \textbf{Alice:} if Bob plays ST\#1, ST\#1 (10) $>$ ST\#2 (4). If Bob plays ST\#2, ST\#1 (8) $>$ ST\#2 (5). \textbf{ST\#1 is strictly better in BOTH cases} $\Rightarrow$ \textbf{Alice has a strictly dominant strategy: ST\#1.}
  \item \textbf{Bob:} if Alice plays ST\#1, ST\#1 (10) $>$ ST\#2 (4). If Alice plays ST\#2, ST\#1 (8) $>$ ST\#2 (5). \textbf{ST\#1 is strictly better in BOTH cases} $\Rightarrow$ \textbf{Bob also has a strictly dominant strategy: ST\#1.}
\end{itemize}
\textit{This is the cleanest possible outcome: both players have a dominant strategy, both dominant strategies coincide, and the resulting cell is the unique NE --- a ``dominant strategy equilibrium''.}
\end{solbox}

%% ==========================================
\section{2023 --- Q4(b): Design a game-theory strategy to identify attackers in a mobile social network (3)}
%% ==========================================

\begin{qbox}{Question --- 2023 Q4(b)}
Using game theory, design a strategy to identify attackers in a mobile social network. \hfill [3]
\end{qbox}

\begin{solbox}{Answer --- 2023 Q4(b) (Lc\#3 ``Identifying Attackers in a Mobile Social Network'')}
\textbf{Model the network as a game between two players:}
\begin{itemize}[noitemsep]
  \item \textbf{Player 1 --- the Server} (network/security monitor), which connects with each node.
  \item \textbf{Player 2 --- a Node}, which can be either \textbf{benign} (a normal user) or \textbf{malicious} (an attacker) --- the server does not directly know which.
\end{itemize}

\textbf{Actions available to each player:}
\begin{itemize}[noitemsep]
  \item \textbf{Server's actions:} (1) Do nothing, (2) Send the node a normal packet, (3) Set up surveillance on the node.
  \item \textbf{Node's actions:} (1) Do nothing / ignore, (2) Forward the packet normally, (3) Damage / drop the packet (malicious behavior).
\end{itemize}

\textbf{Payoff structure (the strategy design):}
\begin{itemize}[noitemsep]
  \item If the server \textbf{sends a normal packet}: good outcome if the node is \textbf{benign} (forwards it normally), bad outcome if the node is \textbf{malicious} (damages it).
  \item If the server \textbf{surveils}: good outcome if the node is \textbf{malicious} (caught/detected), but imposes a cost/bad outcome if the node was actually \textbf{benign} (wasted resources, degraded service).
  \item If the server \textbf{does nothing}: safe but no information gained --- a malicious node could infiltrate undetected.
\end{itemize}

\textbf{The strategy:} the server should choose a \textbf{mixed strategy} --- i.e., randomize between ``send normal traffic'' and ``surveil'' with some probability for each node, rather than a fixed rule. This is because:
\begin{enumerate}[noitemsep]
  \item Always surveilling everyone $\Rightarrow$ degraded service for all benign (majority) nodes.
  \item Never surveilling $\Rightarrow$ malicious nodes infiltrate freely.
  \item By analyzing the \textbf{payoff matrix} of (server action $\times$ node type/action) and finding the \textbf{Nash equilibrium} of this game, the server can compute the \textbf{optimal surveillance probability} that balances:
  \begin{itemize}[noitemsep]
    \item maximizing detection of malicious nodes (security), against
    \item minimizing disruption to benign nodes (quality of service).
  \end{itemize}
\end{enumerate}

$\Rightarrow$ \textbf{Conclusion:} game theory provides a mathematical framework to automate this trade-off analysis at scale (\emph{``analyze hundreds of thousands of `what ifs'\,''}), giving network administrators a principled, equilibrium-based surveillance policy instead of an ad-hoc one.
\end{solbox}

%% ==========================================
\section*{2020--2022 --- No Game Theory question found}
%% ==========================================
\begin{warnbox}{2020/2021/2022 papers}
No Game Theory / Nash Equilibrium / payoff-matrix question was identified in the 2020, 2021, or 2022 past papers. This is a topic that appears to have been newly introduced into the syllabus for 2023+2024 (both most recent years) --- a strong recency signal that it will likely repeat in 2026.
\end{warnbox}

%% ==========================================
\section*{Exam Strategy Summary}
%% ==========================================

\begin{keybox}{High-Yield Checklist for Game Theory (2/5 years, both most recent, 3--6 marks)}
\begin{enumerate}[noitemsep]
  \item \textbf{Memorize the underline method:} for each column, underline the row-player's best (highest) payoff; for each row, underline the column-player's best (highest) payoff. Any cell with BOTH numbers underlined is a Nash Equilibrium --- there can be 0, 1, or 2+.
  \item \textbf{A $2\times2$ matrix can have TWO NEs} (coordination games like 2024 Eve/Mallory) --- don't stop after finding one if it's a "prefer to match" scenario.
  \item \textbf{Strictly dominant strategy} = one action is always best regardless of the opponent's choice. Check each player's payoffs for BOTH of their own actions against EACH of the opponent's actions --- if one row/column choice wins in every comparison, it's dominant.
  \item \textbf{If both players have the same dominant strategy} and it coincides with the unique NE (2023 Alice/Bob, both ST\#1), explicitly say so --- this is the "cleanest" answer and often worth bonus clarity marks.
  \item \textbf{Conceptual variant} (2023 Q4(b)): "design a strategy" $\to$ describe the server/node model (actions, payoffs, surveillance trade-off) and conclude with "find the Nash equilibrium of the payoff matrix to balance security vs. service quality."
  \item \textbf{Always show the labeled payoff matrix with underlines} --- this is the visual proof that earns full method marks even if the final NE identification has a typo.
\end{enumerate}
\end{keybox}

\end{document}
